\datos{color}{}{}
\section*{Tarea}
\addcontentsline{toc}{section}{Tarea}
%\addcontentsline{toc}{section}{Ejercicio 1}
Como uno de los objetivos fundamentales de la asignatura es conseguir que los
alumnos sean capaces de enfrentarse a un problema de optimización real con 
las técnicas aprendidas, el primer trabajo que se propone consiste en que 
cada alumno sugiera el problema de optimización que desee resolver a lo largo
de la asignatura, identifique sus variables de decisión, funciones objetivo,
restricciones, grupo al que pertenecen, etc. \\

Con el objeto de nivelar los problemas elegidos por los alumnos, 
los problemas elegidos serán propuestos por los alumnos a los
profesores de la asignatura, y analizados, modificados, refinados,
y finalmente aceptados o rechazados por los últimos hasta
alcanzar una propuesta formal del problema que los alumnos irán resolviendo
a lo largo de la asignatura. \\

Este proceso de definición y selección de problemas se realizará durante el estudio de los contenidos de este primer tema.
Debido al trabajo en equipo que se establece entre los profesores y los alumnos, se proponen dos fechas finales de entrega de
la definición del problema:

\begin{itemize}
    \item 15 de Marzo: Fecha limite de entrega de la propuesta inicial del problema en la que el alumno presentará un documento que
describa (cualitativamente) las características mas relevantes del problema que desea resolver, identificando las variables de
decisión y las funciones de mérito utilizadas para medir los objetivos y restricciones del problema.

    \item Durante las dos semanas sucesivas, los profesores y el alumno analizaran el problema, lo modificarán, etc... con el objeto
de acordar una propuesta inicial definitiva antes del 29 de Marzo. Para esta fecha, el alumno documentará la propuesta
definitiva a los profesores, poniendo de manifiesto las variables de decisión, funciones de merito, restricciones, clasifi-
cación del problema, etc... Además, implementará las funciones responsables de medir la bondad (mérito y restricciones)
de una posible solución del problema. En las sesiones de dudas del tema (Semana 23 - 27 Febrero) se propondrá a
los alumnos la cabecera de la función, de forma que sea utilizada, sin cambios significativos, durante los siguientes temas.
\end{itemize}